% Laget automatisk av elevark/lag_elevark.py fra opplegg/fotoner-fotoelektrisk.html. Endringer her blir overskrevet.
\documentclass[11pt]{article}
\usepackage{elevark}
\begin{document}
\elevarkhode{Fysikk 2 · Kvantefysikk}{Fotoner og fotoelektrisk effekt – arbeidsark}{https://fysikklab.no/sim/fotoelektrisk.html}{Fotoelektrisk effekt}

\section*{Oppgaver}

Oppgave 1–6 bruker Fotoelektrisk effekt\footnote{\url{fysikklab.no/sim/fotoelektrisk.html}}, og oppgave 7 bruker Dobbeltspalte med én partikkel om gangen\footnote{\url{fysikklab.no/sim/dobbeltspalte.html}}. Oppgave 8–9 er uten simulering. Bruk $h = 6{,}63 \cdot 10^{-34}\,\mathrm{Js} = 4{,}14 \cdot 10^{-15}\,\mathrm{eVs}$ og $c = 3{,}00 \cdot 10^{8}\,\mathrm{m/s}$.

\begin{oppgaveliste}
\oppgave \fase{gjett}{Gjett}Velg sink og rødt lys (650~nm). Skru intensiteten helt opp. Kommer det elektroner? Hva sier bølgemodellen?
\svar
\oppgave \fase{test}{Test}Velg natrium og 400~nm. Endre intensiteten. Hva skjer med strømmen, og hva skjer med den største kinetiske energien?
\svar
\oppgave \fase{test}{Undersøk}Finn den lengste bølgelengden som slår løs elektroner fra natrium. Sammenlign med grensefrekvensen i simuleringen.
\svar
\oppgave \fase{test}{Regn}Regn ut fotonenergien og den største kinetiske energien for natrium med 400~nm. Sjekk med simuleringen.
\svar
\oppgave \fase{test}{Mål}Med natrium og 400~nm: Øk bremsespenningen til strømmen akkurat blir null. Stemmer $eU_{\mathrm{b}}$ med $E_{k,maks}$?
\svar
\oppgave \fase{test}{Mål}Tøm punktene. Lagre punkter for 250, 300, 350, 400 og 450~nm, hver gang med bremsespenningen der strømmen akkurat blir null. Hvilke verdier får du for $h$ og $W$?
\svar
\oppgave \fase{forklar}{Forklar}Send fotoner gjennom dobbeltspalten med 10 per sekund. Hvordan kan hvert foton treffe ett punkt, samtidig som det til slutt blir striper?
\svar
\oppgave \fase{test}{Regn}Hvor mange fotoner per sekund sender en grønn laser på 5,0 mW og 550~nm ut?
\svar
\oppgave \fase{utfordring}{Utfordring}I et forsøk ble bremsespenningen målt til 0,18~V ved $6{,}0 \cdot 10^{14}\,\mathrm{Hz}$, 1,01~V ved $8{,}0 \cdot 10^{14}\,\mathrm{Hz}$ og 1,84~V ved $10{,}0 \cdot 10^{14}\,\mathrm{Hz}$. Finn $h$ og løsrivingsenergien. Hvilket metall kan det være?
\svar
\end{oppgaveliste}

\end{document}
